%% appendix_assumptions.tex — Declared Assumptions and Known Limitations
%% Extracted from main.tex
\chapter{Declared Assumptions and Known Limitations}
\label{app:assumptions}

\begin{enumerate}[leftmargin=*,label=\textbf{A\arabic*.}]

\item \textbf{Conditional feature independence.}  The eight features are
treated as conditionally independent given habitability $H$.  In
practice, organic abundance and chemical energy are positively
correlated, as are liquid hydrocarbon and surface--atmosphere
interaction.  The independence assumption is conservative: correlated
positive features slightly overstate the posterior, but the framework
remains physically motivated and the data-derived importances are
consistent with the prior weights (Section~\ref{sec:importances}).

\item \textbf{Static geomorphological template.}  The Lopes et al.\
(2019) map is used for all five temporal configurations.  The
Late Heavy Bombardment surface had different dune distributions and
crater density, and the future surface will be submerged.  The present
template is used as the best available proxy, with amplitude
modifications applied through epoch-specific prior means and scale
functions.

\item \textbf{Seamless organic abundance mode.}  The organic abundance
feature uses terrain-class scores globally rather than VIMS spectral
band ratios, because the VIMS mosaic covers only $\sim$50\% of
the globe and direct blending produced a visible seam at the coverage
boundary at 180\textdegree{}W.

\item \textbf{Past-epoch Cassini feature contamination.}  The past-epoch
posterior was computed using present-era Cassini SAR maps, which
encode the present polar lake morphology even at $t = -3.5\gya$.
The north polar Bayesian score in the past anchor is inflated above the
physically expected prior value.  This is mitigated in the multi-epoch
reconstruction by using the analytical Bayesian formula directly for
pre-lake-formation frames (Section~\ref{sec:multi_epoch}) and is
documented in Section~\ref{sec:past_contamination}.

\item \textbf{Multi-epoch temporal scaling accuracy.}  The multi-epoch
reconstruction applies scalar scale functions $s_i(t)$ to the present-epoch
feature maps and re-evaluates the Bayesian posterior formula at each epoch
(Section~\ref{sec:epoch_scaling}).  This is physically motivated but
is not equivalent to independent epoch-specific feature extraction.
The scale functions are calibrated to reproduce known physical transitions
(lake formation, solar warming, eutectic crossing) but have no observational
constraint at intermediate epochs.  Estimated uncertainty:
$\Delta P(H) \lesssim \pm0.05$ at non-anchor epochs.

\item \textbf{SAR coverage gap.}  Approximately 33\% of \titan{}'s
surface lacks SAR coverage.  Features 1, 3, and 8 are degraded in
these regions, where geomorphological and topographic proxies
substitute.

\end{enumerate}
